% Solution: Cabral (2017) Ch. 9 Exercise 9.19 "Secret price cuts". Input by \solitem. Check: check_T3.py (item 25).
\begin{exsol}{Cabral Exercise 9.19}{245}{Advanced}{5}
\asked Two duopolists face demand $h = 1$ (probability $1-\alpha$) or $l = 0$ (probability $\alpha$) at reservation
price $u$, with zero cost, and each firm sees only its own demand. They set $p = u$ until a firm gets zero demand,
then $p = 0$ for $T$ periods, then $p = u$ again. Show (a) the value $V$; (b) the deviation value $V'$; (c) the
equilibrium condition; (d) the optimal equilibrium.

\method Cabral \S9.2 (pp.~222--227; secret price cuts, Green--Porter). The strategy is an equilibrium iff the value of
cooperation $V$ is at least the value $V'$ of the best deviation (one-stage deviation principle). The price war
itself is the static equilibrium ($p = 0$), so it needs no check.

\solution
\textbf{(a) $V$.} In a cooperative period, with probability $1-\alpha$ demand is high and each firm sells
$\frac12$ at $u$, earning $\frac u2$ and staying in the cooperative phase (value $\delta V$ next period). With
probability $\alpha$ demand is low: profit $0$, and the price war follows for $T$ periods of zero profit, after which
cooperation returns, a value $\delta^{T+1}V$ viewed from today. Hence
\[
  V = (1-\alpha)\Bigl(\frac u2 + \delta V\Bigr) + \alpha\,\delta^{T+1}V
  \quad\Longrightarrow\quad
  V = \frac{(1-\alpha)u/2}{1 - (1-\alpha)\delta - \alpha\delta^{T+1}} .
\]
\textbf{(b) $V'$.} A deviator undercuts $u$ slightly. If demand is high (probability $1-\alpha$) he sells $1$ unit at about
$u$, earning $u$, and the rival, who then gets zero demand, starts the price war. If demand is low (probability $\alpha$)
he earns $0$ either way and the war starts anyway. In both cases the future is $\delta^{T+1}V$ (war and return):
\[
  V' = (1-\alpha)u + \delta^{T+1}V .
\]
\textbf{(c) Condition.} $V \ge V'$ means $V\,(1 - \delta^{T+1}) \ge (1-\alpha)u$. From (a), $(1-\alpha)u =
2V\,[\,1 - (1-\alpha)\delta - \alpha\delta^{T+1}\,]$. Substituting and dividing by $V > 0$:
\[
  1 - \delta^{T+1} \ge 2 - 2(1-\alpha)\delta - 2\alpha\delta^{T+1}
  \iff 1 \le 2(1-\alpha)\delta + (2\alpha - 1)\delta^{T+1} .
\]
\textbf{(d) Optimal equilibrium.} The planner chooses $T$ to maximise $V$, which is increasing in $\delta^{T+1}$, so a
\emph{shorter} war is better, subject to (c). Write $g(T) = 2(1-\alpha)\delta + (2\alpha - 1)\delta^{T+1}$.
\begin{itemize}
  \item If $\alpha \ge \frac12$: $g(T) \le \delta < 1$ for every $T \ge 0$: no such equilibrium exists. (Wars are
  triggered too often.)
  \item If $\alpha < \frac12$: $g$ increases in $T$ and tends to $2(1-\alpha)\delta$, so an equilibrium exists iff
  $\delta > \frac{1}{2(1-\alpha)}$, and the optimal one uses the \emph{smallest} $T$ with
  $\delta^{T+1} \le \dfrac{2(1-\alpha)\delta - 1}{1 - 2\alpha}$.
\end{itemize}
Example: $\alpha = \frac15$, $\delta = \frac9{10}$: $g(1) = 0.954 < 1 \le g(2) = 1.0026$, so $T^* = 2$ and $V = 2.98$ (with $u = 1$).
Equilibrium price wars occur on the path (after bad draws) even though nobody cheats: they are the cost of
imperfect monitoring.

\interp A price war must be long enough to deter secret cuts and short enough not to waste surplus. The optimal
war is the shortest one that deters, and it happens with positive probability on the equilibrium path.

\verify \texttt{check\_T3.py} (item~25) solves (a) symbolically, derives $(V - V')$ times the denominator and
shows it equals $g(T) - 1$ times $\frac{(1-\alpha)u}{2}$ for $T = 1,2,3$, confirms $V = 2.98$ by value iteration over the
Markov chain of phases ($T = 2$, $\alpha = 0.2$, $\delta = 0.9$), shows that (c) is never met for $\alpha \ge \frac12$,
and finds $T^* = 2$ by search together with $\lceil\ln r / \ln\delta\rceil - 1$ for $r = \frac{2(1-\alpha)\delta - 1}{1-2\alpha}$.
Cabral prints no answer.

\mistakes
\begin{itemize}
  \item Treating the war as a punishment for cheating only: it also follows low demand, which is why the equilibrium
  has wars on the path.
  \item Maximising $T$: a longer war lowers $V$ (more periods of zero profit).
\end{itemize}

\answer{(a) $V = \frac{(1-\alpha)u/2}{1-(1-\alpha)\delta-\alpha\delta^{T+1}}$. (b) $V' = (1-\alpha)u + \delta^{T+1}V$.
(c) $V \ge V'$ is equivalent to $1 \le 2(1-\alpha)\delta + (2\alpha-1)\delta^{T+1}$. (d) Exists iff $\alpha < \frac12$
and $\delta > \frac1{2(1-\alpha)}$; the optimal war has the smallest $T$ satisfying (c), e.g.\ $T^* = 2$ for $\alpha =
\frac15$, $\delta = 0.9$.}
\end{exsol}
